\section{Actual quotient costs controlled by the edge zero potential}
\label{sec:edge-potential-quotient-cost}

This section estimates the canonical cost of MA-PR-010 in its
original whole-quotient norm.  It gives sufficient local zero
geometry for polynomial costs, and records exactly what the
standard count and generic distance estimates alone supply.
It asserts neither that the new sufficient geometry holds
for the fixed arithmetic function nor that a contrary
arithmetic zero configuration occurs.

Keep the fixed primitive unitary Hecke function, pole-removed
entire completion \(\Xi\), Hilbert space and actual quotients
of (MA-PR.1).  Write
\[
 e_-=a+u,\qquad e_+=a+v,\qquad \nu_n=2+|n|,
 \qquad
 \Gamma_n=\partial\bigl([-1,2]+i[h_n,h_{n+1}]\bigr),
 \tag{MA-PR.150}
\]
with the original selected heights and holomorphic quotient
norms \(J_n\) of MT-M18.  Assume the generic edge-distance
condition in exactly its proved form:
\[
 \min\{|\operatorname{Re}\rho-e_-|,
               |\operatorname{Re}\rho-e_+|\}
       \ge c_0(2+|\operatorname{Im}\rho|)^{-2}
                     \qquad(\Xi(\rho)=0).
 \tag{MA-PR.151}
\]
All zeros are counted with their actual multiplicities.
There are no boundary zeros, so every retained primary
factor is retained in full, for every fixed \(s\ge0\).
The canonical maps and costs are therefore defined:
\[
 \mathcal C_n
   =\|\iota_n\mathscr R_n\|_{\mathcal L(Q_H,J_n)}.
 \tag{MA-PR.152}
\]
The norm on the source is the norm on the whole \(Q_H\);
in particular no oblique finite-block projection is
discarded or replaced by an orthogonal one.

Let
\[
 \begin{split}
 \mathcal Z_n&=\{\rho:n-2\le\operatorname{Im}\rho\le n+3\},\\
 U_n&=\max_{e\in\{e_-,e_+\}}
       \sup_{h_n\le t\le h_{n+1}}
       \sum_{\rho\in\mathcal Z_n}m_\rho
             \log^+\frac1{|e+it-\rho|}.
 \end{split}
 \tag{MA-PR.153}
\]
This is a finite nonnegative real number.  Unlike a
bound using only real-part distances, it retains the
actual distances in the complex plane and the spread
of zero ordinates within a block.

\paragraph{MA-PR-012 (whole-quotient cost and a sufficient local count).}
Under (MA-PR.151), there are fixed finite constants
\(C,D\), depending on the fixed arithmetic function
and Hilbert parameters, such that
\[
                    \mathcal C_n\le C\nu_n^D e^{U_n}
                                   \qquad(n\in\mathbb Z).
 \tag{MA-PR.154}
\]
Thus the sufficient condition \(U_n=O(\log\nu_n)\)
gives all the actual returns, coherent lifts and
all-order nuclear representations of MA-PR-010
for every \(g\in\mathcal E\).
This conclusion holds for every fixed \(s\ge0\).

Here are two explicit sufficient conditions for that
potential bound.  Either may be imposed at both edges.
\begin{enumerate}
\item There are constants \(A_0,A_1\) such that, for every
block \(n\), \(t\in[h_n,h_{n+1}]\), and \(0<r\le1\),
\[
 N_e(t,r):=\sum_{|e+it-\rho|<r}m_\rho
                   \le A_0+A_1r\log\nu_n .
 \tag{MA-PR.155}
\]
\item There are fixed \(\eta>0\) and \(M<\infty\) such that
\[
 \sum_{\substack{\rho\in\mathcal Z_n\\
                     |\operatorname{Re}\rho-e|<\eta}}
                         m_\rho\le M
                         \qquad(n\in\mathbb Z).
 \tag{MA-PR.156}
\]
\end{enumerate}
The first condition permits a growing total block count
provided it has the stated distribution at small radii.
The second is a bounded count only in a fixed band about
the edges, rather than in the whole critical strip.
Neither condition is deduced here from the ordinary
\(O(\log\nu_n)\) count.

Without either additional condition, the standard count
and (MA-PR.151) yield only the established estimate
\[
                U_n=O(\log^2\nu_n),\qquad
                \mathcal C_n\le C\exp(C'\log^2\nu_n).
 \tag{MA-PR.157}
\]
This is an upper bound supplied by these estimates,
not a lower bound or an asserted arithmetic failure of
polynomial costs.  Absence of proper boundary truncation
does not improve this argument by itself.

\begin{proof}
\emph{1. The actual inner restriction and the full output class.}
Choose a fixed small \(\delta_0>0\) and set
\[
 \delta_n=\delta_0\nu_n^{-2},\qquad
 \Omega_n^{\rm in}
  =\{e_-+\delta_n<\operatorname{Re}w<e_+-\delta_n,\
                       h_n<\operatorname{Im}w<h_{n+1}\}.
 \tag{MA-PR.158}
\]
Take \(\delta_0<(v-u)/4\).  For \(\rho\in\mathcal Z_n\),
\(2+|\operatorname{Im}\rho|\) and \(\nu_n\) are comparable
with absolute constants.  Decrease \(\delta_0\) so that
\[
 \delta_n\le\tfrac14
           \min\{|\operatorname{Re}\rho-e_-|,
                          |\operatorname{Re}\rho-e_+|\}
                      \qquad(\rho\in\mathcal Z_n).
 \tag{MA-PR.159}
\]
The rectangle contains exactly the retained zeros of
block \(n\), with their full multiplicities.  Its height,
width and perimeter have the same uniform bounds as
in (MA-PR.6).  Empty retained blocks have cost zero
and may be omitted.

For \(h\in H\), the function
\(F_h(w)=\mathcal Mh(w-a)\) is holomorphic on the
open Hilbert strip and continuous on this closed
inner rectangle.  The actual weighted
Cauchy--Schwarz estimate gives
\[
 \sup_{\partial\Omega_n^{\rm in}}|F_h|
 \le
 \|h\|_H
 \sup_{u+\delta_n\le b\le v-\delta_n}
 \left(\int_{\mathbb R}
      \frac{e^{2bx}}{\max(e^{2ux},e^{2vx})(1+x^2)^{s/2}}
                  \,dx\right)^{1/2}
 \le C\nu_n\|h\|_H .
 \tag{MA-PR.160}
\]
Only \(s\ge0\) is used: discard the polynomial denominator
and integrate \(e^{-2\delta_n|x|}\).
Let \(J_n^{\rm in}\) be the quotient of
\(A(\Omega_n^{\rm in})\) by the full retained jets,
with its boundary supremum quotient norm.  The
restriction in (MA-PR.160), followed by this quotient,
annihilates \(T_a\mathcal I\), hence its closure \(J_H\).
It is consequently a continuous map from the whole
\(Q_H\), of norm at most \(C\nu_n\).

Fix \(\tau>0\).  For a representative \(F\) of a class
in \(J_n^{\rm in}\), use the actual entire interpolant
\[
 (\mathcal H_nF)(z)
   =\frac1{2\pi i}\int_{\partial\Omega_n^{\rm in}}
       F(w)\frac{\Xi(z)-\Xi(w)}{\Xi(w)(z-w)}
                        e^{\tau(z-w)^2}\,dw .
 \tag{MA-PR.161}
\]
The divided difference is entire in \(z\), including
\(z=w\).  The same Cauchy argument as (MA-PR.9)
gives \(\mathcal H_nF=F+\Xi A\) in the inner
rectangle, where \(A\) is holomorphic there, and
divisibility by \(\Xi\) near every exterior zero.
If the full retained jets of \(F\) vanish, then
\(F/\Xi\) is holomorphic inside; both exterior
Cauchy integrals vanish.  Entireness makes
\(\mathcal H_nF=0\) everywhere.  Thus (MA-PR.161)
depends on the actual quotient class.

For each fixed \(n\), its entire function decreases
rapidly on every vertical strip: the fixed-strip
bound for \(\Xi\), its derivative bound near the
removable diagonal, and the Gaussian factor give
this directly.  Mellin inversion as in MT-M18
therefore supplies an actual element of \(\mathcal E\).
It has the input's complete retained jets and zero
full jets at every discarded zero and every other
block.  In particular, on taking its class in the
\emph{original full} \(J_n\), one obtains exactly
\[
 \bigl[\mathcal H_nF_h\bigr]_{J_n}
                 =\iota_n\mathscr R_n(\pi h).
 \tag{MA-PR.162}
\]
This identity follows from all primary Taylor classes,
not from a trace comparison.  It proves which actual
whole-quotient operator the following norm estimate
measures.

\emph{2. The potential controls the genuine contour denominator.}
For \(e=e_-\) or \(e_+\), let \(e_n\) be its displaced
inner edge.  Equation (MA-PR.159) gives, for every
\(\rho\in\mathcal Z_n\),
\[
 |e_n+it-\rho|\ge\tfrac34|e+it-\rho|.
 \tag{MA-PR.163}
\]
Indeed the real component changes by at most one
quarter of its absolute value and the imaginary
component does not change.  Thus the sum of
logarithmic losses on either inner edge is at most
\(U_n+(\log(4/3))\sum_{\mathcal Z_n}m_\rho\).
Zeros outside \(\mathcal Z_n\) have imaginary
distance greater than one from every such edge
point: \(h_n>n+1/4\) and \(h_{n+1}<n+7/4\).
They contribute no positive logarithm of an inverse
complex distance.

At large height use the fixed finite disk centered
at \(2+it\), with radius chosen so that both fixed
edge real parts and their small displacements lie
in its inner quarter.  The radius ranges over a
fixed interval and its boundary avoids zeros.
There are \(O(\log\nu_n)\) zeros in the disk.
The Euler lower bound at its center, polynomial
finite-\(L\) upper bound, and Harnack argument of
MT-M18 give a polynomial inverse bound for the
zero-free factor.  A normalized Blaschke factor
satisfies
\[
 |B_\rho(w)|\ge (2R)^{-1}\min\{1,|w-\rho|\}.
 \tag{MA-PR.164}
\]
All factors with distance at least one cost only
\((2R)^{-1}\); their product is a fixed power of
\(\nu_n\).  The remaining factors are controlled
by (MA-PR.163) and \(U_n\).
Multiplying by the lower Stirling bound for the
completion therefore gives, on the two inner
vertical edges,
\[
 |\Xi(w)|\ge c\nu_n^{-D_0}
                         e^{-\kappa|n|}e^{-U_n},
 \qquad \kappa=\pi[K:\mathbb Q]/4.
 \tag{MA-PR.165}
\]
On the horizontal edges the selected-height
bound (M18.8) is polynomial with the same
completion exponential.  Enlarging \(D_0\)
gives (MA-PR.165) on the entire inner contour.
The finitely many small heights have nonzero
contour denominators, so adjustment of the
constant covers them too.

\emph{3. Estimate in the original full block norm.}
For \(z\in\Gamma_n\) and
\(w\in\partial\Omega_n^{\rm in}\), all real parts
range over one fixed interval, and their imaginary
parts differ by less than \(3/2\).  The derivative
bound following (M18.6), applied on the segment
between \(z\) and \(w\), gives
\[
 \left|\frac{\Xi(z)-\Xi(w)}{z-w}\right|
             \le C\nu_n^{D_1}e^{-\kappa|n|}.
 \tag{MA-PR.166}
\]
This includes the removable value at \(z=w\).
The Gaussian factor is bounded on these pairs,
and the inner perimeter is bounded.  Equations
(MA-PR.161), (MA-PR.165) and (MA-PR.166) yield
\[
 \sup_{z\in\Gamma_n}|\mathcal H_nF(z)|
       \le C\nu_n^{D_2}e^{U_n}
                        \sup_{\partial\Omega_n^{\rm in}}|F|.
 \tag{MA-PR.167}
\]
Taking the infimum over representatives gives
the same bound from \(J_n^{\rm in}\) to the
original \(J_n\).  Equivalently, (MA-PR.167)
directly applies to \(F_h\); (MA-PR.160) then
bounds the operator on \(H\) by
\(C\nu_n^{D_2+1}e^{U_n}\).
Its kernel contains \(J_H\), and (MA-PR.162)
identifies its quotient map.  Infimizing over
the representatives of \(\pi h\) proves (MA-PR.154).
This is a bound on \(\iota_n\mathscr R_n\) on
all of \(Q_H\).  No separate estimate on an
oblique \(P_n\), or scalar interpolation surrogate,
has replaced that norm.

\emph{4. The stated zero-count conditions and the exact loss.}
For \(t\in[h_n,h_{n+1}]\), every zero in
\(|e+it-\rho|<1\) belongs to \(\mathcal Z_n\).
The layer-cake identity for this finite multiset is
\[
 \sum_{\rho\in\mathcal Z_n}m_\rho
          \log^+\frac1{|e+it-\rho|}
                =\int_0^1N_e(t,r)\,\frac{dr}{r}.
 \tag{MA-PR.168}
\]
Equation (MA-PR.151) makes this count zero when
\(r<\varepsilon_n:=\min(1,c_*\nu_n^{-2})\), for a fixed
\(c_*>0\).  Under (MA-PR.155), the integral is at most
\[
 A_0\log(1/\varepsilon_n)+A_1(1-\varepsilon_n)\log\nu_n
                                      =O(\log\nu_n).
 \tag{MA-PR.169}
\]
Under (MA-PR.156), the at most \(M\) near-edge zeros
cost \(M\log(1/\varepsilon_n)\).  Every other zero in the
window costs at most \(\log^+(1/\eta)\); the ordinary
window count makes their sum \(O(\log\nu_n)\).
This again bounds \(U_n\) by \(O(\log\nu_n)\).

If only the ordinary bound
\(N_e(t,r)\le C\log\nu_n\) is available, (MA-PR.168)
instead gives
\[
       \int_{\varepsilon_n}^1 C\log\nu_n\,\frac{dr}{r}
                              =O(\log^2\nu_n).
 \tag{MA-PR.170}
\]
This is the explicit product-of-multiplicities loss;
the whole-cluster construction has not silently
replaced it by a fixed exponent.  Formula (MA-PR.154)
now gives (MA-PR.157).  If the stronger potential
bound holds, it gives the polynomial cost condition
of MA-PR-010, whose complete proof supplies all
the asserted actual returns and topology.
\end{proof}

\paragraph{A separate exact single-jet lower bound.}
The potential is sufficient, not asserted necessary.
There is nevertheless an intrinsic norm lower bound
which is useful when assessing possible improvements.
For this separate statement, (MA-PR.151) may be
replaced by the no-proper-truncation assumption
(MA-PR.100) on every primary factor, so that
(MA-PR.152) remains defined.
Let \(\rho=\beta+i\gamma\) be any fully retained zero
in block \(n\), and \(0\le j<m_\rho\).  Put
\[
 d_\Omega(\rho)
      =\min\{1,\operatorname{dist}(\rho,\Gamma_n)\}.
\]
Using the actual norms of (MA-PR.152), one has
\[
 \mathcal C_n\ge
 \frac{d_\Omega(\rho)^j}{j!}
 \left(\int_{\mathbb R}
       \frac{x^{2j}e^{2(\beta-a)x}}{w_{u,v,s}(x)}\,dx\right)^{1/2}.
 \tag{MA-PR.171}
\]
The hypothesis that this factor is fully retained
ensures that every displayed jet is continuous.
This statement also applies when a fully retained
boundary factor exists, although the main upper
bound above was proved with the generic hypothesis.

To prove (MA-PR.171), Cauchy's estimate on disks
of radius increasing to \(d_\Omega(\rho)\) gives
\[
        \left|\frac{F^{(j)}(\rho)}{j!}\right|
               \le d_\Omega(\rho)^{-j}\|[F]\|_n .
 \tag{MA-PR.172}
\]
This is well defined on the full holomorphic quotient.
Its composite with \(\iota_n\mathscr R_n\) is exactly
the retained Hilbert jet in (MA-PR.46).
That functional annihilates \(J_H\), so its norm
on \(Q_H\) equals its weighted Hilbert dual norm;
the square of that norm is the integral in
(MA-PR.171) divided by \(j!^2\).
Taking the supremum over the whole quotient unit
ball proves the lower bound.

The original selected-height product bound also gives
fixed \(c,D_3>0\) such that
\[
              d_\Omega(\rho)^j\ge c\nu_n^{-D_3}
                         \qquad(0\le j<m_\rho).
 \tag{MA-PR.173}
\]
Indeed each horizontal distance, truncated at one
and raised to \(m_\rho\), is at least the entire
product in (M18.7).  Both adjacent products are
bounded below by a fixed power of \(\nu_n\).
The vertical boundaries are at least one from
the critical strip.  Hence
\(d_\Omega(\rho)^{m_\rho}\ge c\nu_n^{-D_3}\);
since \(0<d_\Omega\le1\) and \(j<m_\rho\), this
implies (MA-PR.173).

For a strictly interior zero at right-edge distance
\(\delta=e_+-\beta>0\), and \(2j-s+1>0\), restricting
the integral to \(x\ge1\) gives the more explicit
bound
\[
 \mathcal C_n\ge
 \frac{2^{-s/4}d_\Omega(\rho)^j}{j!}\,
 (2\delta)^{-j+(s-1)/2}
                  \Gamma(2j-s+1,2\delta)^{1/2},
 \tag{MA-PR.174}
\]
where \(\Gamma(q,z)=\int_z^\infty t^{q-1}e^{-t}\,dt\).
This follows from
\((1+x^2)^{s/2}\le2^{s/2}x^s\) for \(x\ge1\).
The left edge has the identical bound with
\(\delta=\beta-e_-\), by using \(x\le-1\).
These inequalities do not assert any growing
multiplicity or very small distance for actual
arithmetic zeros.  They identify an actual
whole-quotient jet norm which cannot be deleted
merely by replacing separate interpolation
formulas by a whole-cluster formula.

\paragraph{Remaining arithmetic scope and source placement.}
The ordinary count and the generic distance are
proved arithmetic inputs; their consequence here
is (MA-PR.157).  The stronger local disk or edge-band
conditions are additional explicitly stated sufficient
conditions.  Their validity at a particular arithmetic
width, or polynomial costs without them, is not
established.  In particular this argument does not
close the unrestricted PR011 gate from absence of
proper truncation alone.  It does not prove a
full common-space matrix theorem, replace the
boundary distributions of MA-MO-027/030, or remove
the discarded-divisor trace correction.

The inputs are the complete MT-M18 block theorem,
selected-height and minimum-modulus arguments,
the actual inner restriction and contour construction
MA-PR-001, the whole-quotient Gram and module
identifications MA-PR-006, and the proved
polynomial-cost implication MA-PR-010.
The new estimates keep those actual norms throughout.
Their finite-length module and nuclear-trace meaning
is the one already specified in (MA-PR.80) and
MA-PR-009--011.  No new reading of the inherited
cote~1 scans, additional outside arithmetic premise,
or Grothendieck-group splitting is asserted.
